\section{Sufficiency and Safety Metrics (ESI, CR@k, RPI@k)}
\label{sec:c1_metrics}

\subsection{Metric overview}
Component 1 computes four metric families from a claim-level partition $(\Active,\Susp,\Counter)$:
\begin{itemize}
\item Evidence Sufficiency Index: $\ESIprod$, $\ESIgeom$.
\item Counter Retention: $\CR@k$.
\item Recovery Potential Index: $\RPI@k$.
\item Auxiliary uncertainty: neutral dominance in Active set.
\end{itemize}

\begin{figure}[H]
\centering
\begin{adjustbox}{max width=\textwidth}
\begin{tikzpicture}[
  node distance=7mm and 9mm,
  block/.style={draw,rounded corners,align=center,text width=3.25cm,minimum height=1.15cm,fill=gray!7},
  arrow/.style={-Latex,thick}
]
\node[block] (A) {Active set $\Active$};
\node[block, right=of A] (mass) {Mass\\$\hat{M}_{\Active}$};
\node[block, above right=of mass] (cov) {Coverage\\$\Cov_{\Active}$};
\node[block, below right=of mass] (conn) {Connectivity\\$\Conn_{\Active}$};
\node[block, right=24mm of mass] (esi) {ESI outputs\\$\ESIprod,\;\ESIgeom$};

\draw[arrow] (A) -- (mass);
\draw[arrow] (A) |- (cov);
\draw[arrow] (A) |- (conn);
\draw[arrow] (mass) -- (esi);
\draw[arrow] (cov) -- (esi);
\draw[arrow] (conn) -- (esi);
\end{tikzpicture}
\end{adjustbox}
\caption{Page-fitted ESI computation structure from Active evidence.}
\label{fig:c1_esi_flow}
\end{figure}

\subsection{ESI components as implemented in \texttt{compute\_esi}}

\paragraph{Mass term.}
\begin{align}
M_{\Active} &= \sum_{e\in\Active} \Rel(e), \\
\alpha &= \max(1, |\Active|), \\
\hat{M}_{\Active} &= 1-\exp\!\left(-\frac{M_{\Active}}{\alpha}\right).
\end{align}

\paragraph{Coverage and connectivity raw terms.}
Let $E_{\text{claim}}$ be claim entities and $V(\Active)$ graph nodes from Active triples.
\begin{align}
\Cov_{\Active,\text{raw}} &= \frac{|E_{\text{claim}}\cap V(\Active)|}{|E_{\text{claim}}|},
\end{align}
with edge-case override to 1.0 if $|E_{\text{claim}}|=0$.

Connectivity is computed over present claim entities:
\begin{equation}
E_{\text{present}} = E_{\text{claim}}\cap V(\Active).
\end{equation}
If $|E_{\text{present}}|<2$, connectivity raw is set to 1.0.
Otherwise:
\begin{equation}
\Conn_{\Active,\text{raw}} = \frac{\#\{\{u,v\}\subseteq E_{\text{present}}: u\leftrightsquigarrow v\}}{\binom{|E_{\text{present}}|}{2}}.
\end{equation}
Connectivity uses BFS on a simple undirected graph built from subject-object edges.

\paragraph{Smoothing and final ESI forms.}
With $\epsilon=0.05$:
\begin{align}
\Cov_{\Active} &= \epsilon + (1-\epsilon)\Cov_{\Active,\text{raw}}, \\
\Conn_{\Active} &= \epsilon + (1-\epsilon)\Conn_{\Active,\text{raw}}.
\end{align}
Final scores:
\begin{align}
\ESIprod &= \hat{M}_{\Active}\cdot\Cov_{\Active}\cdot\Conn_{\Active}, \\
\ESIgeom &= \left(\hat{M}_{\Active}\cdot\Cov_{\Active}\cdot\Conn_{\Active}\right)^{1/3}.
\end{align}

\paragraph{Why both are logged.}
\begin{itemize}
\item $\ESIprod$ is strict and can collapse under conservative NLI outputs.
\item $\ESIgeom$ is less collapse-prone and used as primary metric in summaries/examples.
\end{itemize}

\subsection{Counter Retention at $k$}
\texttt{compute\_cr\_at\_k} uses keyword-only signature:
\begin{lstlisting}[style=jsonstyle]
compute_cr_at_k(*, pool_items, counter_items, ks=[5, 10])
\end{lstlisting}
This prevents silent positional argument swapping.

Let \texttt{TopContra\_k(Pool)} be top-$k$ by $\Con(e)$:
\begin{equation}
\CR@k = \frac{|\texttt{TopContra\_k}(\Pool)\cap\Counter|}{\min(k,|\Pool|)}.
\end{equation}
The denominator fix \(\min(k,|\Pool|)\) is critical for small pools.

\subsection{Recovery Potential Index at $k$}
\texttt{compute\_rpi\_at\_k} estimates connectivity gain from Suspended bridges.

\paragraph{Bridge definition in code.}
A suspended triple $e=(u,r,v)\in\Susp$ is a bridge candidate iff:
\begin{enumerate}
\item both endpoints $u,v$ exist in Active graph nodes,
\item $u$ and $v$ belong to different connected components of Active graph.
\end{enumerate}
Candidates are sorted by $\Rel(e)$ descending.

For each $k\in\{1,3,5\}$ by default:
\begin{align}
\Active^{(+k)} &= \Active \cup \text{top-}k\text{ bridges from }\Susp, \\
\RPI@k &= \Conn(\Active^{(+k)}) - \Conn(\Active).
\end{align}

\subsection{Neutral dominance}
\begin{equation}
\text{NeutralRate}_{\Active} =
\begin{cases}
0, & |\Active|=0,\\
\frac{1}{|\Active|}\sum_{e\in\Active} \Neu(e), & \text{otherwise.}
\end{cases}
\end{equation}
High neutral rate often indicates verbalization uncertainty or off-topic evidence.

\subsection{Interpreting metrics on large runs (real train summary)}
From \texttt{logs/component1/train/summary.json}:
\begin{table}[H]
\centering
\begin{tabular}{@{}p{0.42\linewidth}p{0.20\linewidth}p{0.30\linewidth}@{}}
\toprule
Statistic & Value & Interpretation \\
\midrule
\texttt{mean\_esi\_geom} & 0.3151 & Moderate average sufficiency. \\
\texttt{esi\_geom\_p10} & 0.1000 & Bottom-decile starvation reference. \\
\texttt{starve\_rate\_geom\_fixed\_03} & 0.6358 & 63.58\% claims below 0.3 threshold. \\
\texttt{mean\_coverage\_A} & 0.7796 & Entities frequently present but not always complete. \\
\texttt{mean\_connectivity\_A} & 0.3961 & Structural disconnection is common. \\
\texttt{mean\_neutral\_rate\_A} & 0.6113 & PV often remains neutral-dominant. \\
\bottomrule
\end{tabular}
\caption{Selected real summary metrics (train split).}
\label{tab:c1_real_summary_stats}
\end{table}
